% !TeX root = ../navier-stokes-manuscript.tex
% ============================================================
% 2.1 Critical Height at a Finite Endpoint
% ============================================================

\subsection{Critical Height at a Finite Endpoint}\label{sub:CriticalHeightAtAFiniteEndpoint}

The one-fluid picture leaves a Blown failure one place to hide: the field can remain smooth at every earlier time while its activity is driven into smaller scales before a finite endpoint.
Fix a rapidly decaying datum \(u_0\) on \(\R^3\), a viscosity \(\nu>0\), and the maximal classical solution \((u,p)\) they generate on \([0,T_*)\).
Suppose \(T_*<\infty\) so that we can follow that threat all the way to its alleged endpoint.

Kinetic energy stays bounded along the classical interval, but a fixed amount of energy is not equally visible at every scale.
The Navier--Stokes scaling shows the difference.
For every \(\lambda>0\), the rescaled fields
\[
  u_\lambda(x,t)
  :=
  \lambda u(\lambda x,\lambda^2t),
  \qquad
  p_\lambda(x,t)
  :=
  \lambda^2p(\lambda x,\lambda^2t)
\]
solve the same equation with the same viscosity on the rescaled time interval.
Length has been divided by \(\lambda\), time by \(\lambda^2\), and velocity has been multiplied by \(\lambda\).
The two spatial derivatives in viscosity and the one time derivative in acceleration are already setting the pace.

Let \(\Lambda=(-\Delta)^{1/2}\), and measure a spatial Sobolev height by
\[
  \norm{v}_{\dot H^s}^2
  :=
  \int_{\R^3}
  |\xi|^{2s}|\widehat v(\xi)|^2
  \,d\xi.
\]
The Fourier transform of the rescaled velocity is
\[
  \widehat{u_\lambda}(\xi,t)
  =
  \lambda^{-2}
  \widehat u(\xi/\lambda,\lambda^2t).
\]
Putting this into the Sobolev norm and changing variables from \(\xi\) to \(\eta=\xi/\lambda\) gives
\begin{align*}
  \norm{u_\lambda(t)}_{\dot H^s}^2
  &={}
  \int_{\R^3}
  |\xi|^{2s}\lambda^{-4}
  \left|\widehat u(\xi/\lambda,\lambda^2t)\right|^2
  \,d\xi
  \\
  &={}
  \lambda^{2s-1}
  \int_{\R^3}
  |\eta|^{2s}
  \left|\widehat u(\eta,\lambda^2t)\right|^2
  \,d\eta.
\end{align*}
Hence
\[
  \norm{u_\lambda(t)}_{\dot H^s}
  =
  \lambda^{s-1/2}
  \norm{u(\lambda^2t)}_{\dot H^s}.
\]
Every Sobolev height tilts under the compression except \(s=1/2\).
At that height the power of \(\lambda\) vanishes, so define
\[
  H(t)
  :=
  \frac12\norm{u(t)}_{\dot H^{1/2}}^2
  =
  \frac12\norm{\Lambda^{1/2}u(t)}_2^2.
\]
This is the critical height: exact rescaling can move a smooth configuration to a finer length and a shorter time without changing the value of \(H\).

The kinetic-energy scale moves in the other direction.
Putting \(s=0\) in the same calculation gives
\[
  \norm{u_\lambda(t)}_2^2
  =
  \lambda^{-1}
  \norm{u(\lambda^2t)}_2^2.
\]
A finer rescaled structure can carry less \(L^2\) energy while standing at the same critical height.
That is the opening left by the ordinary energy estimate: it can keep the total energy finite while the scale at which the fluid is changing continues to fall.

Finite time makes that opening sharper.
Choose successively smaller lengths
\[
  \ell_k
  :=
  2^{-k}\ell_0.
\]
The parabolic time attached to the \(k\)-th length has size
\[
  \tau_k
  :=
  \frac{c\ell_k^2}{\nu}
  =
  \frac{c\ell_0^2}{\nu}4^{-k},
\]
where \(c>0\) fixes the normalization.
These times fit inside a finite interval because
\[
  \sum_{k=0}^{\infty}\tau_k
  =
  \frac{4c\ell_0^2}{3\nu}
  <
  \infty.
\]
The sum makes the Zeno picture dimensionally possible: smaller and smaller parabolic episodes could be scheduled before one finite time while their critical height remains visible.
It does not show that the Navier--Stokes evolution generates those episodes.
The actual fluid would still have to create each finer state from the state before it while transport, pressure, incompressibility, and viscosity all continue to act.

Critical height sees what can survive the loss of scale.
It still says nothing about how rapidly the field must change to pass through infinitely many regular states before \(T_*\).
Finite time turns that missing speed into a relation between one response and the response immediately above it.
